\documentclass[12pt,a4paper]{article}
\usepackage[utf8]{vietnam}
\usepackage{amsmath,amssymb}
\usepackage{geometry}
\geometry{
  top=12mm,
  bottom=12mm,
  left=14mm,
  right=14mm,
  footskip=6mm,
  headheight=14pt
}

\usepackage{tikz}
\usetikzlibrary{arrows.meta,calc,shapes.geometric}
\usepackage[most]{tcolorbox}
\usepackage{fontawesome}
\usepackage{fancyhdr}
\usepackage{enumitem}

\pagestyle{fancy}
\fancyhf{}
\fancyhead[L]{\footnotesize\color{teal!80!black}\faBook\ \textbf{SỔ TAY KIẾN THỨC TOÁN 3}}
\fancyhead[R]{\footnotesize\color{teal!80!black}\textbf{Lớp Toán Thầy Lĩnh -- loptoanthaylinh.vn}}
\fancyfoot[C]{\footnotesize\color{gray}-- Sổ tay kiến thức: Bài 05. Ôn tập bảng chia 2, bảng chia 5 --}
\renewcommand{\headrulewidth}{0.4pt}
\renewcommand{\footrulewidth}{0pt}

\setlength{\parindent}{0pt}
\setlength{\parskip}{0pt}

\begin{document}

% --- BANNER BÀI HỌC ---
\begin{tcolorbox}[
  enhanced,
  colback=teal!8!white,
  colframe=teal!70!black,
  arc=3mm, boxrule=1.2pt,
  left=3mm, right=3mm, top=2mm, bottom=2mm,
  halign=center
]
  {\LARGE\bfseries\color{teal!90!black}\faBookmark\ BÀI 05: ÔN TẬP BẢNG CHIA 2, BẢNG CHIA 5}\\[1mm]
  {\itshape\footnotesize\color{black!75} Mối quan hệ phép nhân và phép chia, kĩ thuật hình thành bảng chia 2 và 5 trên 1 trang A4}
\end{tcolorbox}

\vspace{-2mm}

\noindent
\begin{minipage}[t]{0.49\textwidth}\vspace{0pt}
\begin{tcolorbox}[
  enhanced,
  colback=blue!4!white,
  colframe=blue!75!black,
  arc=2mm, boxrule=0.8pt,
  left=2mm, right=2mm, top=1.8mm, bottom=1.8mm,
  title={\footnotesize\bfseries\color{white}\faShareAlt\ 1. Bảng chia 2 (Chia đều cho 2)}
]
{\footnotesize
\begin{itemize}\setlength{\itemsep}{1pt}\setlength{\leftskip}{-2.5mm}
  \item \textbf{Từ phép nhân suy ra phép chia:}
  $$2 \times k = P \implies P : 2 = k$$
  \item Chia cho 2 nghĩa là tìm \textbf{một nửa} của một số.
  \item Số chia hết cho 2 phải là các số chẵn tận cùng $0, 2, 4, 6, 8$.
  \item \textit{Ví dụ:} $14 : 2 = 7$; $18 : 2 = 9$; $20 : 2 = 10$.
\end{itemize}
}
\end{tcolorbox}
\end{minipage}
\hfill
\begin{minipage}[t]{0.49\textwidth}\vspace{0pt}
\begin{tcolorbox}[
  enhanced,
  colback=teal!4!white,
  colframe=teal!75!black,
  arc=2mm, boxrule=0.8pt,
  left=2mm, right=2mm, top=1.8mm, bottom=1.8mm,
  title={\footnotesize\bfseries\color{white}\faShareAlt\ 2. Bảng chia 5 (Chia đều cho 5)}
]
{\footnotesize
\begin{itemize}\setlength{\itemsep}{1pt}\setlength{\leftskip}{-2.5mm}
  \item \textbf{Từ phép nhân suy ra phép chia:}
  $$5 \times k = Q \implies Q : 5 = k$$
  \item Số bị chia trong bảng chia 5 luôn có tận cùng là $0$ hoặc $5$.
  \item \textbf{Ý nghĩa thực tế:} Chia thành 5 phần bằng nhau, hoặc chia mỗi nhóm 5 phần tử.
  \item \textit{Ví dụ:} $35 : 5 = 7$; $45 : 5 = 9$; $50 : 5 = 10$.
\end{itemize}
}
\end{tcolorbox}
\end{minipage}

\vspace{-2.5mm}
% --- KHỐI 3: SƠ ĐỒ TƯ DUY TIKZ ---
\noindent
\begin{tcolorbox}[
  enhanced,
  colback=orange!4!white,
  colframe=orange!80!black,
  arc=2mm, boxrule=0.8pt,
  left=2mm, right=2mm, top=1.5mm, bottom=1.5mm,
  title={\footnotesize\bfseries\color{white}\faSitemap\ SƠ ĐỒ TƯ DUY: CẦU NỐI 2 CHIỀU PHÉP NHÂN \& PHÉP CHIA}
]
\centering
\begin{tikzpicture}[
  box/.style={draw, rounded corners=2pt, align=center, font=\scriptsize\bfseries, inner sep=3pt, minimum height=0.85cm},
  arr/.style={-{Stealth[length=2.5mm]}, thick}
]
  \node[box, fill=purple!12, draw=purple!80!black, text width=3.3cm] (m1) at (0, 0) {PHÉP NHÂN\\$2 \times 8 = 16$\\$5 \times 7 = 35$};
  \node[box, fill=teal!12, draw=teal!80!black, text width=3.3cm] (m2) at (5.1, 0) {CHIA CHO THỪA SỐ 1\\$16 : 2 = 8$\\$35 : 5 = 7$};
  \node[box, fill=blue!12, draw=blue!80!black, text width=3.3cm] (m3) at (10.2, 0) {CHIA CHO THỪA SỐ 2\\$16 : 8 = 2$\\$35 : 7 = 5$};

  \draw[arr, purple!80!black] (m1) -- (m2);
  \draw[arr, teal!80!black] (m2) -- (m3);
\end{tikzpicture}
\end{tcolorbox}

\vspace{-2.5mm}
% --- KHỐI 4: VÍ DỤ & MẸO TRÁNH BẪY ---
\noindent
\begin{minipage}[t]{0.49\textwidth}\vspace{0pt}
\begin{tcolorbox}[
  enhanced,
  colback=green!4!white,
  colframe=green!60!black,
  arc=2mm, boxrule=0.8pt,
  left=2mm, right=2mm, top=1.8mm, bottom=1.8mm,
  title={\footnotesize\bfseries\color{white}\faCheckCircle\ Ví dụ thực hành chuẩn mực}
]
{\footnotesize
\textbf{1. Tính nhẩm nhanh:}\\
$12 : 2 = \mathbf{6}$;\quad $18 : 2 = \mathbf{9}$.\\
$25 : 5 = \mathbf{5}$;\quad $40 : 5 = \mathbf{8}$.\\
\textbf{2. Tìm một nửa:}\\
Một nửa của $16$ quả cam là: $16 : 2 = \mathbf{8}$ (quả).\\
\textbf{3. Chia phần:}\\
Có $30$ cái kẹo chia đều cho $5$ bạn. Mỗi bạn được: $30 : 5 = \mathbf{6}$ (cái kẹo).
}
\end{tcolorbox}
\end{minipage}
\hfill
\begin{minipage}[t]{0.49\textwidth}\vspace{0pt}
\begin{tcolorbox}[
  enhanced,
  colback=red!4!white,
  colframe=red!70!black,
  arc=2mm, boxrule=0.8pt,
  left=2mm, right=2mm, top=1.8mm, bottom=1.8mm,
  title={\footnotesize\bfseries\color{white}\faExclamationTriangle\ Mẹo tránh bẫy trắc nghiệm}
]
{\footnotesize
\begin{itemize}\setlength{\itemsep}{1.5pt}\setlength{\leftskip}{-2.5mm}
  \item \textbf{Mẹo nhẩm phép chia bằng phép nhân:}\\
  Muốn tính $35 : 5 = ?$, hãy nhẩm: "Mấy nhân với $5$ bằng $35$?". Vì $5 \times 7 = 35$ nên kết quả là $7$.
  \item \textbf{Bẫy chia cho 0:}\\
  Số chia luôn phải khác $0$. Không bao giờ có phép chia cho $0$!
  \item \textbf{Phân biệt:}\\
  $10 : 2 = 5$ (chia đều cho 2 bạn).\\
  $10 - 2 = 8$ (bớt đi 2 cái).
\end{itemize}
}
\end{tcolorbox}
\end{minipage}


\end{document}
